\chapter{Benchmarking Methodology}
\label{ch:benchmarking-methodology}

\Cref{ch:usecases-design} fixed what is asked of a detector. This chapter fixes how
the answer is measured. It sets the rules that make a comparison fair, the measures
on the three axes of detection quality, computational cost and usability, the
protocols for batch and for streaming evaluation, and the tests that turn per-seed
numbers into statements. It is the methodological contribution C1 and the basis of
the answer to RQ5. It reports no result; the numbers are in \cref{ch:results}.

The protocol exists to test nine hypotheses, listed in \cref{tab:hypotheses}. Each
sits under the research question it bears on. A hypothesis is a claim that the data
can contradict, and it is stated here without any conclusion. The set was fixed
before the final runs. The tests of this chapter decide how each is read, and where
the datasets cannot test one, the chapter says so.
% FIXME(open): the wording of H1 to H9 is reconstructed from the meeting of 14 October 2026; check it against the written version of the co-supervisor when it arrives, and decide whether Ch. 1, Sec. 1.3 cites \cref{tab:hypotheses} in one sentence.

\begin{table}[t]
\centering
\footnotesize
\caption{The hypotheses of the benchmark, by research question, with the experiments
that test them. The knowledge effect is LTN-full minus LTN-sup; the effect of the loss
is LTN-sup minus MLP-CE.}
\label{tab:hypotheses}
\setlength{\tabcolsep}{4pt}
\begin{tabular}{@{}p{1.9cm}p{1.5cm}p{7.2cm}p{2.2cm}@{}}
\toprule
Hypothesis & RQ & Statement & Experiments \\
\midrule
H1 Quality & RQ1
  & A knowledge arm with well-built, tuned axioms reaches a higher macro-F1 than a
    neural network trained without them, on the same data & E2, E5, E6 \\
H2 Rare classes & RQ1, RQ2
  & The gain of the knowledge arms is concentrated in the rare classes & E2, E3, E5, E8 \\
H3 Changed axioms & RQ2
  & If the initial axioms are changed, the gain shrinks, stays or reverses; it
    reverses when LTN-full falls below LTN-sup & E3, E11 \\
H4 Data quantity & RQ2
  & The gain grows as the labelled data shrink and fades as the training set grows & E7 \\
H5 Shift & RQ2
  & The gain changes under a shift of distribution, on attack types unseen in
    training and across datasets & E6, E9 \\
H6 Degradation & RQ2
  & When the axioms are degraded in controlled steps, the performance falls
    gradually and not catastrophically & E11 \\
H7 Cost & RQ3
  & Training cost grows linearly with the number of axioms; at inference an
    LTN-trained network costs about the same as the same network trained with
    cross-entropy & E13, E12 \\
H8 Constraints & RQ4
  & At equal quality, a knowledge arm violates fewer domain rules than a
    data-driven one & E2, E5, E6, E12 \\
H9 Classical ML & RQ5
  & Random Forest and gradient boosting remain strong baselines on tabular flow
    data, at a much lower cost of training & E2, E5, E6, E13 \\
\bottomrule
\end{tabular}
\end{table}

Three readings of the table are fixed now. H1 is tested against both baselines:
LTN-full against MLP-CE, which is the claim as stated, and against LTN-sup, which says
whether the gain comes from the knowledge or from the loss. H2 cannot be tested on
CIC-DDoS2019, which has two classes and no rare one (\cref{sec:uc2}). H5 also asks
whether a longer training buys a better transfer; this is read from the training times
that every run records.

% ================================================================
% 8.1 Fair benchmark
% ================================================================
\section{Requirements for a fair benchmark}
\label{sec:fair-benchmark}

A comparison is fair when the difference between two arms can be attributed to what
changes between them. Five requirements follow. They answer the pitfalls that
\citet{arp2022dosdonts} list for learning-based security systems, and the sources of
variance that \citet{bouthillier2021variance} study in machine-learning benchmarks,
among them the data split, the initialisation and the choice of hyperparameters.
Model evaluation and selection in general are surveyed by \citet{raschka2018model}.
% FIXME(open): \cref{sec:cost-dimension}, \cref{sec:statistical-plan}, \cref{sec:harness} and \cref{sec:threats} render as "??" until Secs. 8.3, 8.6, 8.7 and 8.8 are written.

\subsection{Matched capacity: same architecture, same features, same preprocessing}
\label{sec:matched-capacity}

The four neural arms share one network: two hidden layers of 256 and 128 units with
ReLU, dropout 0.2, one output per class, trained with Adam at batch size 256 and the
selection rule of \cref{sec:matched-tuning}. The arms differ in the loss and in what
they do about class imbalance, which are the things under test. The number of
parameters is recorded in every run, so equality can be checked.
% FIXME(open): the common network is fixed for all three datasets by the tuning files (D6 of the protocol); on the two CIC datasets it is a proposal to confirm, since it was chosen with NSL-KDD in view. It is not optimal for every arm, which \cref{sec:threats} must say.

All arms receive the same features and the same preprocessing. Numeric columns are
standardised and categorical columns one-hot encoded, with the scaler and the encoder
fitted on the training part only. Random Forest and XGBoost read the same matrix. The
axioms are written on the raw columns (\cref{sec:batch-pipeline}), which the knowledge
arms and the rule-only arm read beside the network input. If a feature computed over a
window of flows is added to a dataset, it goes to every arm; otherwise a gain could
come from information that only one arm sees.

Capacity is matched for the four neural arms only. Random Forest and XGBoost have
other kinds of capacity, set by their own tuning with the same budget, and they enter
as the strong tabular baselines of H9, to be compared on quality and on cost. The
rule-only arm has no capacity to match.

Imbalance is not handled the same way, and this is declared and not hidden. MLP-CE
weights the loss by class, and in its weighted variant the exponent of the weights is
tuned. The LTN arms aggregate the supervised axioms per class and draw batches in which
every class has a minimum quota, and the strength of the balancing is tuned. The forest
and the boosted trees use class weights, with a tuned exponent. These are different
levers. Each is searched inside the same budget, and \cref{sec:threats} states what
follows from the difference.

\subsection{Matched tuning budget: same search space and same number of trials}
\label{sec:matched-tuning}

Every arm that has hyperparameters gets 20 trials. A trial is one draw from the search
space of the arm, scored by the validation macro-F1 averaged over three tuning seeds,
100, 101 and 102, which are not among the final seeds. The draws come from a Latin
hypercube with a fixed search seed, so a repeated search tries the same configurations.
The trial with the highest mean is kept, and a tie goes to the earlier one. The test
part is never scored while tuning. The chosen configuration is frozen and run on the
ten final seeds. \Cref{tab:tuning-spaces} gives the spaces.

\begin{table}[t]
\centering
\footnotesize
\caption{Search spaces. Each arm is searched with 20 trials of 3 seeds. Ranges are
uniform, log-uniform (log) or a list of values.}
\label{tab:tuning-spaces}
\setlength{\tabcolsep}{4pt}
\begin{tabular}{@{}p{2.3cm}p{10.3cm}@{}}
\toprule
Arm & Parameters searched \\
\midrule
MLP-CE & Learning rate, log $[3\cdot10^{-4}, 3\cdot10^{-3}]$. Weighted variant: also the
  exponent of the class weights, $[0.25, 1]$ \\
LTN-sup & Learning rate as above; exponent of the batch quota per class, $[0, 0.5]$;
  minimum examples per class in a batch, $\{2,4,8,16\}$; exponent $p$ of the
  aggregator, $\{1,2,4,6\}$ \\
LTN-full & The four of LTN-sup; weight $w$ of the domain axioms, $\{0.1, 0.25, 0.5\}$;
  scale of the widths $s$, log $[0.5, 2]$ \\
Random Forest & Trees $\{100,300,500\}$; depth $\{$none$,10,20,30\}$; minimum leaf size
  $\{1,2,5\}$; features per split $\{\sqrt{\cdot},\log_2,0.3\}$; exponent of the class
  weights, $[0, 1]$ \\
XGBoost & Learning rate, log $[0.03, 0.3]$; depth $\{3,\dots,10\}$; minimum child
  weight, log $[1, 10]$; row and column subsample, $[0.6, 1]$ and $[0.5, 1]$; exponent
  of the class weights, $[0, 1]$ \\
Rule-only & None; nothing is learned \\
\bottomrule
\end{tabular}
\end{table}

Early stopping and the choice of epoch follow one rule in every arm. The criterion is
the mean of the validation macro-F1 over the last five epochs, or rounds for XGBoost,
with a patience of ten and the best weights restored. The validation loss cannot be the
common criterion, because cross-entropy and one minus the satisfaction aggregate are
not on the same scale. The moving mean damps a lucky epoch: the validation part holds
few examples of the rare classes, and one example moves the recall of a class such as U2R
by a large step. The mean does not remove the problem, and \cref{sec:threats} lists it.
The neural arms run at most 100 epochs and XGBoost at most 300 rounds. A cap binds
unevenly, since an arm that converges slowly is cut short, so the share of runs that
end at the cap is reported for every arm.
% FIXME(open): the files of the repository give a cap of 100 epochs on NSL-KDD as well; check that the NSL-KDD final runs used 100 and not 50, and correct this paragraph if they did not.

The budget is equal in trials, not in dimensions. The arms offer from one to six
levers, and twenty draws cover a six-dimensional space less densely than a
one-dimensional one. LTN-full and XGBoost are therefore the arms most exposed to
under-tuning. We take equal trials as the definition of an equal budget, because the
levers of an arm are part of what the arm is. On CIC-IDS2017 the tuning trials train
on a quarter of the Normal rows of the training part, to bound the cost of sixty runs
per arm, while the final runs use all rows. A configuration chosen on a thinned set is
then applied to another class balance, and \cref{sec:threats} says so.

The logic has settings of its own, and they belong to the run and not to the knowledge
base file (\cref{sec:kb-format}). The connectives, product and Reichenbach implication,
are fixed (\cref{ch:ltn}). The exponent $p$ is one value per run, used by the
quantifiers and by the aggregation of the axioms into the satisfaction, and it is
tuned; it is not set axiom by axiom. The supervised and the domain axioms are
aggregated separately and combined with the weight $w$ of \cref{tab:tuning-spaces}, so
that the loss is
\[
  1 - \bigl((1-w)\,\mathrm{SatAgg}_{\mathrm{sup}} + w\,\mathrm{SatAgg}_{\mathrm{domain}}\bigr).
\]
This is the second of the two options that \cref{sec:supervised-domain-axioms} leaves
open, and it makes the weight explicit instead of letting the number of axioms in each
group decide it. The widths $s$ of the graded conditions are fixed in the file and
scaled by one factor, which is tuned. Crisp conditions are used only where E4 asks for
them.

\subsection{No test-set leakage: the validation set is never the test set}
\label{sec:no-leakage}

The test part is the one piece of data that no decision may touch. Every decision
that could use it is fixed on training or validation data. The scaler and the encoder
are fitted on the training part. The thresholds $\tau$ and the widths $s$ of the axioms
are fixed on training data, with the statistics of the axiom, support and confidence,
computed there (\cref{sec:kb-elicitation}). The hyperparameters, the epoch and the
operating point (\cref{sec:operating-point}) are chosen on the validation part. The
primary metric was fixed before the first final run. In code, the flag that scores the
test part is false in every configuration used for tuning, and the tuning routine
refuses a base configuration that sets it. The test part is scored once per seed, in
the final run.

The parts are built as follows. NSL-KDD has an official test file. The validation part
is a stratified 15\,\% of the training file, drawn again for each seed. CIC-IDS2017 has
one labelled file, cut into training, validation and test parts in the proportions
60, 20 and 20, stratified by class. The test cut is drawn once with a fixed seed, so
every arm and every seed meets the same test part, and only the cut between training
and validation follows the seed. CIC-DDoS2019 has a training day and a test day, and
the validation part is a stratified 15\,\% of the training day, drawn for each seed. On
the two CIC datasets, rows with the same features and the same label are removed before
any split, so that an identical row cannot sit in the training and in the test part.
In E9 the held-out fine labels are removed from the training and the validation parts
and all their rows are added to the test part.

The axioms are the place where the rule is easiest to break and hardest to see. For
NSL-KDD the thresholds were fixed on the whole training file, so the validation rows of
every seed shaped them. No test row did, but the validation part is not independent of
the axioms, and a validation score of a knowledge arm is slightly flattered by it.
\Cref{sec:threats} records this.

\subsection{Repetition: $n$ independent seeds, report mean and dispersion}
\label{sec:repetition}

Every arm runs on the same ten seeds, 0 to 9. A seed fixes the initial weights, the
order and composition of the batches, the randomness of the trees and the cut between
training and validation. It does not move the test part, which is the same for every
seed. The spread over seeds therefore measures the randomness of training and of the
validation cut. It does not measure the sampling variation of the test set, and it says
nothing about other datasets. The arms share their seeds, and so their training and
validation cuts, which makes the comparison paired: it is made seed by seed
(\cref{sec:statistical-plan}). The three tuning seeds are not among them.

We report, for every measure and arm, the mean and the sample standard deviation over
the seeds, and for a comparison the paired median of the differences with an interval.
Ten is a budget decision, set by the cost of the runs, and not the outcome of a power
analysis. With ten pairs the exact two-sided Wilcoxon test cannot give a $p$-value
below $2/2^{10}\approx 0.002$, a limit that matters when several comparisons share the
correction (\cref{sec:statistical-plan}).

\subsection{Reproducibility: pinned environment, fixed seeds, published artefacts}
\label{sec:reproducibility}

The environment is pinned in a lock file of package versions. A run seeds the Python
hash, the Python and NumPy generators and TensorFlow, and switches on the deterministic
operations of TensorFlow. Its result file records the configuration, the seed, the
versions of Python, TensorFlow, \texttt{ltn}, NumPy, scikit-learn and XGBoost, the
platform, the commit of the code and whether the working tree was modified, the time of
each phase and the metrics of each part. The final runs on the CIC datasets also save
the model and the preprocessing, so that a measure added later can be computed without
training again.
% FIXME(open): check that the NSL-KDD final runs saved the model and the preprocessing too.

Determinism is a property of one machine. On one machine, in the checks we made, two
runs of the same configuration and seed give the same macro-F1 and select the same
epoch, and the loss differs at most in its last digits on the GPU. Across machines and
library versions we do not claim bit-identical results. The final runs were made on a
server with Python 3.14 and TensorFlow 2.22, and the early work on a laptop with Python
3.11 and TensorFlow 2.21. For this reason times are compared only between runs of one
machine (\cref{sec:cost-dimension}), and a metric is not assumed to be equal across the
two.
% FIXME(open): the result files of the NSL-KDD final runs show the server environment, while the plan was to make them on the laptop; confirm which machine made them and correct this paragraph and the cost section.
% FIXME(open): the final runs made on the server record the commit as unknown; record the commit of each series from the repository history (point 22 of the open list) before any figure is called citable.
% FIXME(open): the package sets TF_ENABLE_ONEDNN_OPTS to 0 unless the environment sets it; check that the final runs on both machines used the same value.

What is published is the code, the configurations, the knowledge base files, the
result files and the saved models, in the repository \texttt{fil2405/nesy-nids-bench}.
\Cref{sec:harness} shows how a table is rebuilt from them.

% ================================================================
% 8.2 Detection quality
% ================================================================
\section{Dimension 1 -- Detection quality}
\label{sec:detection-quality-dimension}

Detection quality is measured on the validation part while tuning and on the test part
in the final runs. The measures are chosen for a detector whose Normal class is the
large majority and whose attack classes are not equally rare. Only the macro-F1 is
used to select anything.

\subsection{Metric set and why each one is included}
\label{sec:metric-set}

\Cref{tab:metric-set} lists the measures. The primary one is the macro-F1, the
arithmetic mean of the per-class F1. It is primary because the task has classes of very
different size, and a detector that ignores the rare ones has to pay for it. Accuracy
does not charge that price: a detector that always answers Normal scores high accuracy
on a mostly benign test part and low macro-F1. The macro-F1 was fixed as the primary
measure, and as the criterion of selection in tuning, before the first final run. It is
not changed after the data are seen. Every other measure may qualify a result and none
may replace the primary one.

\begin{table}[t]
\centering
\footnotesize
\caption{Measures of detection quality. All are computed for every arm.}
\label{tab:metric-set}
\setlength{\tabcolsep}{4pt}
\begin{tabular}{@{}p{3.5cm}p{6.4cm}p{3cm}@{}}
\toprule
Measure & What it reads & Role \\
\midrule
Macro-F1 & Mean of the per-class F1, every class with the same weight & Primary \\
MCC & One number from the whole confusion matrix & Secondary \\
Accuracy & Share of correct decisions & Reported; never used to claim \\
Precision, recall, F1 per class & Where the errors are & Secondary; answers ``on which classes'' \\
PR-AUC per class & Ranking of one class against the rest, without a threshold & Secondary; always read beside the macro-F1 \\
FPR at fixed TPR, TPR at fixed FPR & Attack against Normal, on the score $-P(\text{Normal})$ & Secondary; operator view \\
False alarms per 100\,000 benign flows & FPR times 100\,000 & Secondary; alert budget \\
Satisfaction and violation rate per axiom & Agreement of the prediction with the knowledge & Secondary; contribution C5 \\
\bottomrule
\end{tabular}
\end{table}

The Matthews correlation coefficient summarises the confusion matrix in one number that
does not reward the majority class \citep{chicco2020mcc}. It is secondary because it
mixes the classes into one value, where the question of RQ1 is which classes gain.
Precision and recall are therefore reported per class, with the confusion matrix. The
precision--recall area of a class is the average precision of that class against all
the others, from the class probability. The precision--recall view is the more
informative one when a class is rare \citep{davis2006relationship,saito2015precision}.
The operator view reads the attack classes together against Normal, on the score
$-P(\text{Normal})$, with the ROC curve \citep{fawcett2006introduction}. The score is
the negative of the Normal probability and not one minus it, because in single
precision one minus a very small probability rounds to one and creates ties among the
most confident flows. The levels of the operator view are fixed in advance: a
true positive rate of 0.90, 0.95 and 0.99, and a false positive rate of 0.001 and 0.01.
They are starting values of ours, not taken from the literature.

Two conventions are stated. The mean of the macro-F1 runs over all classes of the
dataset, so a class that is absent from both the truth and the prediction of a part
scores an F1 of zero and lowers the mean. This does not happen on the full test parts,
but it can happen in E7 at small fractions of the training data. The precision--recall
area of an absent class is not a number, and the summaries skip it and report how many
seeds they used.
% FIXME(open): decide before E7 whether an absent class keeps the zero F1 or is left out of the mean; the code uses the zero.

\subsection{Threshold-free vs threshold-dependent metrics; choosing an operating point}
\label{sec:operating-point}

The macro-F1, the MCC, the precision and the recall are read at the argmax: the class
with the highest score is the decision. This is what a detector does when nothing else
is tuned, and it is why the macro-F1 is the primary measure.

It has a weakness for this comparison. An arm that balances the classes raises the
recall of the rare classes and lowers their precision. It moves the implicit threshold.
The macro-F1 at the argmax rewards or punishes the move whether or not the arm ranks the
flows better. Two arms can differ in macro-F1 and not in how well they separate the
classes, and without a measure that ignores the threshold the two cases cannot be told
apart. The precision--recall area per class and the operator measures are such
measures. They read the ranking of the scores, and a different threshold does not
change them.

The two readings can diverge, and the protocol does not decide in advance which one
wins. A statement that an arm separates the classes better needs a measure without a
threshold. A statement that it detects better as it would run needs the macro-F1 at the
argmax. When they disagree, both are reported, and the thesis does not pick the one
that is more favourable to an arm.

The operator measures that each run writes have a limit. They are read off the curve of
the part scored, so the threshold is the best one for that very part. They summarise
the ranking, and they are not figures that a deployed detector would reach, because it
cannot choose its threshold on the data it is going to meet. A deployable figure fixes
the operating point on the validation part and applies it unchanged to the test part
(UC2 and UC4 in \cref{ch:usecases-design}). We propose the following rule. The
operating point of an arm is the threshold on $-P(\text{Normal})$ at which the false
positive rate on the validation part equals the alert budget of
\cref{sec:alert-budget}. We report the true positive rate and the false positive rate
that this threshold gives on the test part, and the gap between that false positive
rate and the budget. The recall on an attack unseen in training (FR4) is read at the
same threshold.
% FIXME(open): proposal, to confirm: the operating point at the validation FPR equal to the alert budget (0.001 by default). The transfer of the validation threshold to the test part is not in metrics.py; it needs a script over the saved models and the validation scores of each final run. Ch. 7 says ``operating point fixed on training data'' in FR4 and UC4 and ``on validation data'' in UC2: align Ch. 7 to the validation part.

\subsection{Cost-sensitive evaluation and the alert-budget view}
\label{sec:alert-budget}

A false alarm and a missed attack do not cost the same, and the ratio is unknown. There
is no deployment from which to take it, and an invented ratio would give a score with
no meaning. We therefore build no cost-weighted score. Cost curves show the performance
of a classifier over the whole range of misclassification costs and class priors
\citep{drummond2006cost}, but they need a range that we could not justify.

What can be argued is the base rate. Attacks are rare in real traffic, so a small false
positive rate still produces many alerts that are not attacks, and the analyst stops
trusting the detector \citep{axelsson2000baserate}. The alert budget makes this
operational. It is the number of false alarms that an analyst can handle in a period.
We express it in false alarms per 100\,000 benign flows, which is the false positive
rate times 100\,000. We do not count total alerts per 100\,000 flows, because that
figure depends on the share of attacks in the test part, and the test parts differ
widely in it. The question the budget asks is then this: what true positive rate does
an arm reach at the false positive rate that the budget allows?

On the two CIC datasets the budget is also given per day, from the number of benign
flows per day of the full capture. NSL-KDD has no time stamps, so it has no daily
figure. The daily figure describes a day of a laboratory network and is not a forecast
for a real one. The resolution of the levels depends on the test part. On KDDTest+,
which has 9\,711 benign records, the smallest non-zero false positive rate is about
$10^{-4}$, so a level below that cannot be resolved.
% FIXME(open): the benign flows per day of CIC-IDS2017 and CIC-DDoS2019 are given in Ch. 9, \cref{ch:implementation}, Sec. 9.2.

\subsection{Rule-violation rate: definition and rationale}
\label{sec:rule-violation-rate}

This section defines contribution C5. Take an axiom $j$ with target class $c_j$. It is
evaluated crisp: its antecedent holds on a flow $x$ when every condition holds at its
threshold $\tau$, which is the value at which the graded degree of a condition is $0.5$.
The covered set $C_j$ is the set of flows on which the antecedent holds, and the
coverage is $|C_j|/n$. The \emph{domain-rule violation rate} of the axiom is
\[
  v_j = \frac{\bigl|\{x\in C_j : \arg\max_c P(c\mid x) \neq c_j\}\bigr|}{|C_j|}.
\]
A flow violates the axiom when the antecedent holds and the predicted class is not the
target. This closes the question that \cref{sec:querying} leaves open: the threshold that
turns a degree into a violation is $0.5$ on the antecedent, and the decision is the
argmax. The \emph{axiom-level satisfaction} is
\[
  s_j = 1 - \sqrt{\frac{1}{|C_j|}\sum_{x\in C_j}\bigl(1 - P(c_j\mid x)\bigr)^2},
\]
the $p$-mean-error aggregate with $p=2$ of the Reichenbach implication on the covered
flows, which for a crisp antecedent is the probability of the target class. The \emph{rule precision}
$r_j$ is the share of covered flows whose true class is $c_j$. Both rates are
undefined when no flow is covered, and the summaries then skip the axiom and report the
number of seeds used. The mean violation rate of a model is the unweighted mean over
the axioms that cover at least one flow. Every arm is scored, whether it was trained on
the axioms or not, and always against the original file of the knowledge base, also when
a run trained on an edited copy (E3, E11, E13), so the numbers stay comparable between
arms and between edits.

The evaluation is crisp for a reason. A graded antecedent depends on the width $s$,
which is tuned for each LTN run and does not exist for a forest or a boosted tree.
Crisp conditions give a number that reads as an analyst reads the rule, and that every
arm can be scored on.

Satisfaction and violation are different measures, and neither replaces the other.
Satisfaction uses the probabilities. It separates a target that wins the argmax with a
probability of 0.6 from one that wins with 0.99, and the violation rate does not. With
$p=2$ the aggregate also hides the distribution: very different situations give the
same number, as shown in \cref{sec:querying}. The violation rate is a count that is
easy to read and, for a human, closer to a rule that was broken. The two can rank two
models differently, and both are reported.

The violation is defined against the rule and not against the label, and this limits
how it is read. A rule that is wrong on a flow is violated by the correct answer. A low
violation rate is therefore not a merit in itself: a model that copies the rules has
none, and the rule-only arm has none except where two rules conflict. The rate is read
beside the macro-F1 and the rule precision. For an axiom $j$ the violation rate of a
classifier with perfect labels is $1-r_j$, which is the floor of that axiom. H8 asks for
fewer violations at an equal quality, and it never asks for fewer violations alone.

C5 extends the satisfaction that \citet{onchis2022neurosymbolic} report for rules on
the labels to rules on the features, one rule at a time and for every arm. The
violation rate and the rule precision are our own additions.
